% ===== PRÉAMBULE CANONIQUE — à coller VERBATIM en tête de CHAQUE .tex =====
\documentclass[11pt,a4paper]{article}
\usepackage[utf8]{inputenc}\usepackage[T1]{fontenc}\usepackage{lmodern}
\usepackage[french]{babel}
\usepackage{amsmath,amssymb,mathtools}\usepackage{icomma}
\usepackage[version=4]{mhchem}          % \ce{...} : équations & formules
\usepackage{siunitx}\sisetup{locale=FR} % \qty{}{}, \si{}, \ang{}
\usepackage{chemfig}                    % \chemfig{...} : structures développées
\usepackage{xcolor}\usepackage{colortbl}
\usepackage{tikz}\usepackage{pgfplots}\pgfplotsset{compat=1.18}
\usepackage[most]{tcolorbox}\usepackage{enumitem}\usepackage{array,booktabs}
\usepackage[a4paper,margin=2.2cm]{geometry}
\usetikzlibrary{arrows.meta,calc,angles,quotes,positioning,patterns,backgrounds,shapes.geometric}
\definecolor{primary}{RGB}{222,49,99}\definecolor{primarydark}{RGB}{150,22,55}
\definecolor{defcol}{RGB}{0,114,178}\definecolor{methcol}{RGB}{52,92,148}
\definecolor{excol}{RGB}{0,135,90}\definecolor{attncol}{RGB}{200,40,55}
\tcbset{boxbase/.style={enhanced,breakable,boxrule=0.9pt,arc=2.5pt,
  left=3mm,right=3mm,top=2.3mm,bottom=2.3mm,fonttitle=\bfseries,coltitle=white}}
% --- boîtes TUTORIEL (noms EXACTS, ne pas renommer) ---
\newtcolorbox{cle}[1][]{boxbase,colback=defcol!6,colframe=defcol,
  title={\textsc{À retenir}\ifx\relax#1\relax\relax\else\textmd{ : \itshape #1}\fi}}
\newtcolorbox{methode}[1][]{boxbase,colback=methcol!7,colframe=methcol,sharp corners,
  title={\textsc{Méthode}\ifx\relax#1\relax\relax\else\textmd{ --- \itshape #1}\fi}}
\newtcolorbox{exemplebox}[1][]{boxbase,colback=excol!5,colframe=excol,
  title={\textsc{Exemple}\ifx\relax#1\relax\relax\else\textmd{ : \itshape #1}\fi}}
\newtcolorbox{piege}[1][]{boxbase,colback=attncol!6,colframe=attncol,
  title={\textsc{Piège}\ifx\relax#1\relax\relax\else\textmd{ : \itshape #1}\fi}}
\newtcolorbox{remarque}[1][]{boxbase,colback=black!4,colframe=black!45,
  title={\textsc{Remarque}\ifx\relax#1\relax\relax\else\textmd{ : \itshape #1}\fi}}
% --- boîtes EXERCICE / CORRIGÉ (noms EXACTS, auto-comptées) ---
\newtcolorbox[auto counter]{exo}[1][]{boxbase,colback=primary!4,colframe=primary,
  title={\textsc{Exercice}~\thetcbcounter\ifx\relax#1\relax\relax\else\textmd{ --- \itshape #1}\fi}}
\newtcolorbox[auto counter]{corrige}[1][]{boxbase,colback=excol!4,colframe=excol,
  title={\textsc{Corrigé}~\thetcbcounter\ifx\relax#1\relax\relax\else\textmd{ --- \itshape #1}\fi}}
\newcommand{\R}{\mathbb{R}}\newcommand{\N}{\mathbb{N}}\newcommand{\Z}{\mathbb{Z}}\newcommand{\ds}{\displaystyle}
% (un \usepackage{fancyhdr}+en-tête est possible ; il est ignoré côté web)
% =========================================================================
\begin{document}
% THEME_TITLE: La charge formelle
\newcommand{\CF}{\ensuremath{\mathrm{C.F.}}}
\begin{center}
{\Large\bfseries\color{primarydark} Thème 3 — La charge formelle}
\end{center}

\smallskip
Une fois la structure dessinée, on attribue à chaque atome une \textbf{charge formelle} (\CF) :
un outil de comptabilité qui dit si l'atome a, dans cette structure, \emph{plus} ou \emph{moins}
d'électrons que lorsqu'il est neutre et isolé. C'est le juge de paix entre structures concurrentes.

\begin{cle}[Formule de la charge formelle]
\[ \boxed{\;\CF = N_{v}^{\text{atome}} - \Big(n_{\text{libres}} + \tfrac{1}{2}\,n_{\text{liants}}\Big)\;} \]
En pratique, on compte plus vite avec :
\[ \boxed{\;\CF = (\text{électrons de valence}) - 2\times(\text{doublets libres}) - (\text{nombre de liaisons})\;} \]
car chaque liaison n'apporte qu'\emph{un} électron « propre » à l'atome (l'autre est au partenaire).
\end{cle}

\begin{methode}[Calculer une charge formelle]
\begin{enumerate}[leftmargin=1.9em,topsep=2pt,itemsep=1.5pt]
  \item Repérer $N_v^{\text{atome}}$ (numéro de colonne : $\ce{H}=1,\ \ce{C}=4,\ \ce{N}=5,\ \ce{O}=6,\ \ce{Cl}=7$).
  \item Compter les \textbf{doublets libres} de l'atome, les multiplier par $2$.
  \item Compter le \textbf{nombre de liaisons} de l'atome (simple $=1$, double $=2$, triple $=3$).
  \item $\CF = N_v - 2\times(\text{doublets libres}) - (\text{liaisons})$.
\end{enumerate}
\end{methode}

\begin{cle}[Deux propriétés à connaître]
\begin{itemize}[leftmargin=1.5em,topsep=2pt,itemsep=2pt]
  \item La \textbf{somme} des charges formelles $=$ la \textbf{charge globale} de l'édifice
        ($0$ pour une molécule neutre, $q$ pour un ion).
  \item La \textbf{meilleure structure} est celle dont les \CF{} sont \textbf{minimales} (proches
        de $0$), les charges négatives se plaçant sur les atomes les plus \textbf{électronégatifs}.
\end{itemize}
\end{cle}

\begin{exemplebox}[charges formelles dans \ce{CO3^2-}]
Oxygène \textbf{doublement} lié ($2$ doublets libres, $2$ liaisons) :
$\CF = 6 - 4 - 2 = 0$. \quad
Oxygène \textbf{simplement} lié ($3$ doublets, $1$ liaison) : $\CF = 6 - 6 - 1 = -1$. \quad
Carbone ($0$ doublet, $4$ liaisons) : $\CF = 4 - 0 - 4 = 0$. \;
Somme : $0 + 0 + (-1) + (-1) = -2$ : c'est bien la charge de l'ion.
\end{exemplebox}

\begin{piege}
Ne confonds pas la charge formelle (comptabilité d'électrons dans \emph{une} structure) avec la
charge \emph{réelle} d'un atome ou son nombre d'oxydation. Une molécule \textbf{neutre} peut très
bien porter des charges formelles non nulles qui se compensent (ex. \ce{CO}, \ce{O3}).
\end{piege}

\end{document}
